\chapter{Quality Assessment and Quality of Experience for Multimedia Services}

\textit{Multimedia Communications} Marco Cagnazzo. Part of the slides courtesy of Prof. Federica Battisti. Department of Information Engineering (DEI)

\bigskip\hrule\bigskip

\section{Overview}

This note covers the assessment of multimedia quality and the relationship between quality of service, quality of experience, subjective quality evaluation, objective quality evaluation, modern perceptual metrics, and interactive quality evaluation.

The main topics are:

\begin{itemize}
\item Introduction to \textbf{Multimedia Quality Assessment}.
\item \textbf{QoS} versus \textbf{QoE}.
\item \textbf{Subjective Quality Assessment}.
\item \textbf{Objective Quality Assessment}.
\item Modern perceptual metrics.
\item Interactive quality evaluation.
\end{itemize}

\section{Quality Assessment}

\begin{calloutbox}{red!70!black}{Quality Assessment}
\textbf{Statement:} In digital multimedia services, \textbf{quality assessment} associates a quantitative grade with an image, a video, a sound, or any other possible \textit{stimulus}, that is, a \textbf{signal} perceived by a \textbf{human observer}.

\textbf{Interpretation:} Quality assessment turns the perceived quality of multimedia content into a measurable quantity that can be used when designing, comparing, or optimizing multimedia services.
\end{calloutbox}

Since a stimulus is a signal, quality assessment can be seen as a \textit{functional}: an operator that associates a number with one or more multimedia signals.

\begin{itemize}
\item If the number is obtained by applying an algorithm to the signals, the process is an
\begin{itemize}
\item \textbf{objective quality assessment}.
\end{itemize}
\item If the number is obtained by asking people to evaluate the stimuli, the process is a
\begin{itemize}
\item \textbf{subjective quality assessment}.
\end{itemize}
\end{itemize}

What finally matters is the impact on observers, so subjective assessment remains the reference against which objective methods are judged.

\subsection{Why Quality Assessment Is Needed}

Quality assessment plays a major role in the design of multimedia services because it is the "benefit" term in a cost-benefit analysis: it answers the question, "How good is the quality of the delivered signal?"

The cost depends on the application. For compression and transmission systems, the cost is typically the \textbf{coding rate}, but it can also include:

\begin{itemize}
\item the number of operations, or \textbf{time complexity};
\item the amount of memory required, or \textbf{space complexity};
\item \textbf{energy} consumption;
\item \textbf{latency}.
\end{itemize}
\begin{center}\includegraphics[width=0.600\linewidth]{Pasted_image_20260515104515.png}\end{center}
\textit{Figure 1: Rate-distortion curves for HDR MPEG, HDRV, JPEG HDR fullsamp, JPEG HDR precorrect, and JPEG HDR postcorrect, with SNR in dB plotted against bits per pixel.}

\subsection{Quality Rate Trade-Off}

The trade-off between coding rate and quality is central in multimedia communications.

\begin{itemize}
\item Coding rate and, to some extent, complexity are easy to measure.
\item Quality is much harder to measure because it depends on perception.
\item The course presents different ways to assess image and video quality.
\item Mastering the quality/rate trade-off is necessary to make design decisions in multimedia
\begin{itemize}
\item service systems.
\end{itemize}
\end{itemize}

Typical design questions include:

\begin{itemize}
\item How should two compression methods be compared?
\item How should a streaming service be configured?
\item How much rate is worth spending for a given increase in perceived quality?
\end{itemize}

\subsection{Taxonomy of Quality Assessment Methods}
\begin{center}\includegraphics[width=0.600\linewidth]{Pasted_image_20260515104621.png}\end{center}
\textit{Figure 2: Taxonomy tree for image quality assessment, divided into subjective and objective methods. Subjective methods include single stimulus methods such as ACR/MOS, double stimulus methods with reference and impaired stimuli, and pairwise comparison or A/B comparison. Objective methods include full-reference, reduced-reference, and no-reference metrics.}

The taxonomy separates \textbf{subjective} and \textbf{objective} approaches.

Subjective methods include:

\begin{itemize}
\item \textbf{Single stimulus} methods, including ACR and MOS-based evaluation.
\item \textbf{Double stimulus} methods, where reference and impaired versions are shown.
\item \textbf{Pairwise comparison}, where two versions are compared directly.
\end{itemize}

Objective methods include:

\begin{itemize}
\item \textbf{Full-reference (FR)} metrics, where the original signal is available.
\item \textbf{Reduced-reference (RR)} metrics, where partial side information is available.
\item \textbf{No-reference (NR)} metrics, where no original signal is available.
\end{itemize}

\section{Subjective Quality Evaluation}

\subsection{Subjective Assessment}

Subjective assessment directly measures human judgment.

\begin{itemize}
\item \textbf{Pro:} It is the most reliable way to measure the quality of experience of multimedia services.
\item \textbf{Con:} Subjective tests are usually costly and time-consuming, and require many users to
\begin{itemize}
\item provide reliable results.
\end{itemize}
\end{itemize}

To obtain reliable and reproducible subjective results, the following elements must be defined very carefully:

\begin{itemize}
\item \textbf{Laboratory equipment}, including screen and viewing distance, illumination, and room
\begin{itemize}
\item reverberation characteristics.
\end{itemize}
\item \textbf{Data set}, including the original and processed contents and their distribution.
\item \textbf{Test methodology}, including the rating target, the scale, and the type of stimuli.
\item \textbf{Score processing}, including score normalization, outlier detection, MOS and confidence
\begin{itemize}
\item interval computation, and significance tests.
\end{itemize}
\end{itemize}

\begin{calloutbox}{red!70!black}{Subjective Experiments}
\textbf{Statement:} \textbf{Subjective experiments remain the gold standard for perceptual quality evaluation.}

\textbf{Interpretation:} Objective metrics are useful only if they correlate well with what human observers actually perceive.
\end{calloutbox}

\subsection{Lab Environment}

The laboratory environment affects subjective quality scores. General viewing conditions include:

\begin{itemize}
\item room illumination;
\item chromaticity of the background;
\item peak luminance;
\item monitor contrast ratio;
\item ratio between the luminance of the background behind the picture monitor and the peak luminance
\begin{itemize}
\item of the picture.
\end{itemize}
\end{itemize}

Standardized conditions include:

\begin{itemize}
\item viewing distance;
\item observation angle;
\item room environment color scheme.
\end{itemize}

\subsection{Viewing Distances}

The optimal viewing distance depends on the image system, resolution, aspect ratio, and target horizontal viewing angle.

\begin{center}
\begin{tabular}{|l|l|l|l|l|l|}
\hline
\textbf{Image system} & \textbf{Reference} & \textbf{Aspect ratio} & \textbf{Pixel aspect ratio} & \textbf{Optimal horizontal viewing angle} & \textbf{Optimal viewing distance} \\ \hline
720 x 483 & ITU-R BT.601 & 4:3 & 0.89 & 11 degrees & 7 H \\ \hline
640 x 480 & VGA & 4:3 & 1 & 11 degrees & 7 H \\ \hline
720 x 576 & ITU-R BT.601 & 4:3 & 1.07 & 13 degrees & 6 H \\ \hline
1024 x 768 & XGA & 4:3 & 1 & 17 degrees & 4.5 H \\ \hline
1280 x 720 & ITU-R BT.1543 and BT.1874 & 16:9 & 1 & 21 degrees & 4.8 H \\ \hline
1400 x 1050 & SXGA+ & 4:3 & 1 & 23 degrees & 3.3 H \\ \hline
1920 x 1080 & ITU-R BT.709 & 16:9 & 1 & 31 degrees & 3.2 H \\ \hline
3840 x 2160 & ITU-R BT.2020 & 16:9 & 1 & 58 degrees & 1.6 H \\ \hline
7680 x 4320 & ITU-R BT.2020 & 16:9 & 1 & 96 degrees & 0.8 H \\ \hline
\end{tabular}
\end{center}

Here, $H$ denotes the picture height.

\subsection{Monitor Characteristics}

Using monitors with different characteristics produces different subjective image qualities. Relevant characteristics include:

\begin{itemize}
\item monitor processing;
\item resolution;
\item contrast;
\item brightness;
\item motion interpolation;
\item refresh rate;
\item OLED black levels;
\item local dimming.
\end{itemize}

\subsection{Observers}

Observers may be \textbf{expert} or \textbf{non-expert}, depending on the goals of the assessment.

Typical laboratory studies use between 15 and 40 observers. Before a session, observers should be screened for:

\begin{itemize}
\item \textbf{normal visual acuity}, using a Snellen chart or Landolt chart;
\item \textbf{normal color vision}, using selected color vision charts.
\end{itemize}

Experimenters should describe the characteristics of their assessment panels in as much detail as possible, so that the influence of the panel can be investigated in later studies.

\begin{center}\includegraphics[width=0.645\linewidth]{Pasted_image_20260515104824.png}\end{center}
\textit{Figure 6: Snellen visual acuity chart and Ishihara color vision chart showing the number 12.}

\subsection{Stimuli}

The selection of stimuli depends on the goal of the subjective experiment. Examples include comparing compression algorithms or evaluating the performance of an enhancement method.

For \textbf{video stimuli}, it is useful to compare the relative \textbf{Spatial Information} and \textbf{Temporal Information} of the available sequences. If only a small number of test sequences is used, they should span a large portion of the \textit{spatio-temporal information} space.

\subsection{Spatial Information and Temporal Information}

Spatial Information and Temporal Information help characterize video content.

Intuitively:

\begin{itemize}
\item high $SI$ means textured or detail-rich content;
\item low $SI$ means smooth or cartoon-like content;
\item high $TI$ means sports or action scenes;
\item low $TI$ means news or interview content.
\end{itemize}

\begin{center}
\begin{tabular}{|l|l|l|}
\hline
\textbf{Content} & \textbf{SI} & \textbf{TI} \\ \hline
Football match & Medium & High \\ \hline
Anime & Medium & Low \\ \hline
Nature documentary & High & Medium \\ \hline
Video call & Low & Low \\ \hline
\end{tabular}
\end{center}

\begin{calloutbox}{red!70!black}{Spatial Information}
\textbf{Statement:} \textbf{Spatial Information} ($SI$) measures the amount of spatial detail in a picture. It is usually higher for more spatially complex scenes.

\[
SI = \max_{time}\{\sigma_{space}[Sobel(F_n)]\}
\]

\textbf{Interpretation:} The Sobel filter emphasizes contours and edges. The maximum spatial standard deviation across frames summarizes the highest spatial complexity observed in the sequence.
\end{calloutbox}

The $SI$ computation is based on the Sobel filter, a discrete analogue of the spatial gradient. The output of a Sobel filter is large near object contours and small in flat areas.

For each video frame at time $n$, denoted by $F_n$:

\begin{enumerate}
\item Filter the frame with the Sobel filter, obtaining $Sobel(F_n)$.
\item Compute the standard deviation over the pixels, $\sigma_{space}$.
\item Repeat this for every frame, obtaining a time series of spatial activity.
\item Select the maximum value over time.
\end{enumerate}
\begin{center}\includegraphics[width=0.968\linewidth]{Pasted_image_20260515104905.png}\end{center}
\textit{Figure 7: Pipeline for computing Spatial Information, showing $F_n$ followed by Sobel filtering, variance or standard deviation computation, maximum selection, and $SI$. The visual includes an original grayscale image, a colored Sobel-filtered heat map, and a graph of $\sigma_n^2$ over frames 0 to 250.}

For the example video sequence, $SI = 161.6$.

\begin{calloutbox}{red!70!black}{Temporal Information}
\textbf{Statement:} \textbf{Temporal Information} ($TI$) measures the amount of temporal change in a video sequence. It is usually large for high-motion sequences.

\[
M_n(i,j) = F_n(i,j) - F_{n-1}(i,j)
\]

\[
TI = \max_{time}\{\sigma_{space}[M_n(i,j)]\}
\]

\textbf{Interpretation:} $TI$ measures frame-to-frame luminance differences and summarizes the largest spatial variation of these temporal differences.
\end{calloutbox}

Here, $F_n(i,j)$ is the pixel luminance at the $i^{th}$ row and $j^{th}$ column of the $n^{th}$ frame in time. Diverse motion in adjacent frames produces higher $TI$ values.

\begin{center}\includegraphics[width=0.600\linewidth]{Pasted_image_20260515105231.png}\end{center}
\textit{Figure 8: Motion Difference for a given $n$, shown as a color map with scale from -200 to +200, and a graph of the standard deviation as a function of frame index from 0 to 250. The example reports $TI = 39.20$.}

\subsection{Instructions and Test Session}

Assessors should be carefully introduced to:

\begin{itemize}
\item the method of assessment;
\item the types of impairment or quality factors likely to occur;
\item the grading scale;
\item the sequence and timing of the experiment.
\end{itemize}

Training sequences should demonstrate the range and type of impairments to be assessed. They should use illustrative images other than those used in the actual test, but with comparable sensitivity.

A test session should not last more than half an hour. At the beginning of the first session, about five \textbf{dummy presentations} should be introduced to stabilize the observers' opinion. The data from these presentations must not be considered in the results.

The images and impairments should be presented in a pseudo-random sequence and, preferably, in a different sequence for each session.

\begin{center}\includegraphics[width=0.600\linewidth]{Pasted_image_20260515105459.png}\end{center}
\textit{Figure 9: Temporal scheme of a test session from BT.0500-01-2, including training sequences, stabilizing sequences whose results are not processed, a break for observers' questions, and the main part of the test session.}

The \textbf{learning effect} must also be controlled. Participants may:

\begin{itemize}
\item change their criteria;
\item become more sensitive due to training;
\item become less sensitive due to tiredness.
\end{itemize}

To control this effect:

\begin{itemize}
\item run a training session that covers the range of stimuli, so participants can determine their
\begin{itemize}
\item criteria;
\end{itemize}
\item keep sessions short, for example below 20 minutes or below 40 minutes;
\item pay participants;
\item randomize stimuli as much as possible.
\end{itemize}

\subsection{Subjective Test Methodologies}

\subsubsection{Single-Stimulus or Absolute Category Rating}

In \textbf{single-stimulus (SS)} methods, also called \textbf{Absolute Category Rating (ACR)} methods, a single image or sequence of images is presented, and the assessor provides an index for the entire presentation.

The test material may include:

\begin{itemize}
\item only test sequences;
\item both test sequences and their corresponding reference sequences.
\end{itemize}

In the second case, the reference sequence is presented as a free-standing stimulus and rated like any other test stimulus.
\begin{center}\includegraphics[width=0.806\linewidth]{Pasted_image_20260515105531.png}\end{center}
\textit{Figure 10: Software interface for subjective image quality evaluation with a 0 to 5 rating scale.}

\subsubsection{Absolute Category Rating with Hidden Reference}

In \textbf{ACR-HR}, test sequences are presented one at a time and rated independently on a category scale. The procedure \textbf{must include} a reference version of each test sequence, shown as any other test stimulus. This is called a \textbf{hidden reference} condition.

During data analysis, a \textbf{Differential Mean Opinion Score (DMOS)} is computed between each test sequence and its corresponding hidden reference.

\subsubsection{Double Stimulus Impairment Scale}

In the \textbf{Double Stimulus Impairment Scale (DSIS)} method, the reference and processed sequence/image are shown. Viewers rate degradation on a discrete scale:

\begin{itemize}
\item imperceptible;
\item perceptible but not annoying;
\item fair;
\item poor;
\item bad.
\end{itemize}

\begin{center}\includegraphics[width=0.600\linewidth]{Pasted_image_20260515105605.png}\end{center}
\textit{Figure 11: DSIS timing scheme with Reference Video A for 10 s, pause for 1 s, Test Video B for 10 s, pause for 1 s, Reference Video A} for 10 s, pause for 1 s, Test Video B\textit{ for 10 s, pause for 1 s, and Vote 1 for 5 s, with total duration 49 s and a voting scale from bad to excellent.}

The assessor first sees an unimpaired reference, then the same image impaired. The assessor then votes on the second presentation while keeping the first in mind. Images or sequences are shown in random order and with random impairments covering all required combinations.

\subsubsection{Pairwise Comparison}

In the \textbf{pairwise comparison (PC)} method, two video sequences are shown one after another, and the observer selects the sequence with higher quality.

Example task: for video quality ranking among $n$ sequences, "You will see two video sequences one after another. Select the sequence of higher quality."

Figure 12: Two side-by-side images of a corridor with a fire extinguisher, used as an A/B comparison between an original version and a compressed or modified version.

Results are stored in a \textbf{comparison matrix}:

\[
C =
\begin{bmatrix}
0 & 3 & 1 \\
3 & 0 & 1 \\
1 & 3 & 0
\end{bmatrix}
\quad
\begin{array}{l}
C1 \\
C2 \\
C3
\end{array}
\]

The entry $C_{ij} = n$ means that condition $C_j$ was preferred over $C_i$ exactly $n$ times.

\subsection{Interactive Subjective Quality Evaluation}

The slides include a live subjective quality evaluation.

Instructions:

\begin{itemize}
\item Go to \textbf{wooclap.com} with the code \textbf{MFOZHE}, or scan the QR code.
\item Watch the videos and answer the questions.
\item Focus on perceived visual quality.
\item Vote individually using Wooclap.
\item Do not discuss with neighbors.
\item There are no "correct" answers.
\end{itemize}

The live questions are:

\begin{enumerate}
\item \textbf{Q1: Large Distortions Are Easy to Detect.}
\item 
\item \textbf{Q2: Subjective Quality Is NOT Deterministic.}
\item Which level of quality? 1, 2, 3, 4 or 5? Several people voted for 2, 3, 4. The video was encoded with H.264, PSNR = 35.74 dB.  Is "good" with some distortion still.
\end{enumerate}
Did motion matter more than detail?
\begin{enumerate}
\item \textbf{Q3: Different Artefacts Have Different Perceptual Impact.}
\item The first image is just blurred, the second has just normal artifacts. The PSNRs are almost the same. Some observers prefers blurs, other prefer compression artifacts.
\item \textbf{Q4: Compression Artefacts.}
\item Blocking and motion judder are the most visible.
\item \textbf{Q5/Q6: Frame Rate QoE.}
\item For sport: higher Frame rate is better
\item \textbf{Q7: Perceived Quality Depends on the Viewing Conditions.}
\item Small screen: more tolerance for artifacts instead if 65-inch screens.
\item \textbf{Q8: Human perception is noisy and biased.}
\item They are the same
\end{enumerate}
\begin{center}
\includegraphics[width=0.600\linewidth]{Pasted_image_20260515111804.png}\end{center}

The live activity highlights that:

\begin{itemize}
\item human observers do not always agree;
\item different artefacts produce different annoyance;
\item PSNR alone is insufficient;
\item viewing conditions matter;
\item smooth playback is often preferred;
\item QoE is multidimensional.
\end{itemize}

The central question is whether objective metrics can predict human perception.

\subsection{Processing Subjective Scores}

During subjective experiments:

\begin{itemize}
\item different observers give different scores;
\item some artefacts are considered more annoying;
\item preferences are not always consistent;
\item even identical videos can produce different votes.
\end{itemize}

Therefore, subjective quality assessment requires statistical processing.

\begin{calloutbox}{red!70!black}{Mean Opinion Score}
\textbf{Statement:} Measured scores are considered outcomes of a random process. If $N$ observers rate an image with scores $x_i$, the \textbf{Mean Opinion Score (MOS)} is:

\[
MOS = \frac{1}{N}\sum_{i=1}^{N} x_i
\]

\textbf{Interpretation:} MOS estimates the average perceived quality, but it does not describe disagreement, uncertainty, or statistical reliability by itself.
\end{calloutbox}

\begin{center}\includegraphics[width=0.600\linewidth]{Pasted_image_20260515112817.png}\end{center}
\textit{Figure 21: Histogram for $N = 30$ observers with estimated normal distribution and estimated mean. The x-axis is the score from 0 to 100, and the Gaussian model is $\frac{1}{s\sqrt{2\pi}}e^{-\frac{(x-m)^2}{2s^2}}$.}

Assuming the measured distribution is normal, the mean is estimated as the average of observer ratings.

Average quality is not enough. The analysis also needs:

\begin{itemize}
\item variability;
\item confidence;
\item statistical reliability.
\end{itemize}

\begin{calloutbox}{red!70!black}{Standard Deviation and Standard Error}
\textbf{Statement:} The \textbf{standard deviation} measures observer agreement, while the \textbf{standard error} measures the uncertainty of the estimated MOS.

\[
s = \sqrt{\frac{1}{N}\sum_{i=1}^{N}(x_i - m)^2}
\]

\[
SE = \frac{s}{\sqrt{N}}
\]

\textbf{Interpretation:} More observers do not reduce disagreement among observers, but they do reduce the uncertainty of the MOS estimate.
\end{calloutbox}

Here, $N$ is the sample size and $m$ is the mean.
\begin{center}\includegraphics[width=0.600\linewidth]{Pasted_image_20260515112908.png}\end{center}
\textit{Figure 22: Three histograms for $n = 5$, $n = 50$, and $n = 1000$, showing the population distribution in green and the distribution of the mean in red.}

More observers:

\begin{itemize}
\item do \textbf{not} reduce disagreement, so the same $s$ remains;
\item do reduce uncertainty on the MOS estimate.
\end{itemize}

\begin{calloutbox}{red!70!black}{Confidence Interval}
\textbf{Statement:} Once the uncertainty of the MOS estimate is known, a confidence interval can be computed. If the same experiment is repeated with the same number of observers, in 95\% of cases the average value is expected to fall within the confidence interval.

\[
ci = [-1.96 \cdot SE;\; 1.96 \cdot SE]
\]

\textbf{Interpretation:} The confidence interval expresses how reliable the MOS estimate is, given the number of observers and the variability of their scores.
\end{calloutbox}

\begin{center}\includegraphics[width=0.600\linewidth]{Pasted_image_20260515113323.png}\end{center}
\textit{Figure 23: Three histograms for $n = 5$, $n = 50$, and $n = 1000$, with the confidence region highlighted in red.}

\subsection{Comparison of Subjective Methods}

\begin{center}
\begin{tabular}{|l|l|l|l|l|}
\hline
\textbf{Method} & \textbf{Reference visible?} & \textbf{Speed} & \textbf{Accuracy} & \textbf{Typical use cases} \\ \hline
ACR (Absolute Category Rating) & No & High & Medium & Large-scale streaming evaluation \\ \hline
ACR-HR & Hidden reference & High & High & Compression quality studies \\ \hline
DSIS (Double Stimulus Impairment Scale) & Yes & Medium & Very High & Codec validation \\ \hline
PC (Pairwise Comparison) & Relative comparison & Low & Excellent & Fine perceptual ranking \\ \hline
\end{tabular}
\end{center}

The trade-off is:

\begin{itemize}
\item fast methodologies scale better;
\item more accurate methodologies require more time and attention;
\item no single methodology is optimal in every scenario.
\end{itemize}

\begin{center}\includegraphics[width=0.600\linewidth]{Pasted_image_20260515113344.png}\end{center}
\textit{Figure 24: Decision flowchart for selecting subjective assessment methods. If very accurate comparison is needed, choose Pairwise Comparison. Otherwise, if a reference video is available, choose DSIS or ACR-HR. If fast large-scale testing is needed, choose ACR. Otherwise, use hybrid approaches.}

The choice of methodology depends on:

\begin{itemize}
\item accuracy requirements;
\item testing time;
\item observer fatigue;
\item availability of a reference signal.
\end{itemize}

An important issue is the difference between methods that use explicit references, such as DSIS, and methods that do not use any explicit reference, such as ACR, ACR-HR, and PC. Methods without an explicit reference do not test fidelity.

The \textbf{DSIS} method should be used when testing the fidelity of transmission with respect to the source signal.

\begin{itemize}
\item DSIS has long been a key method specified in ITU-R BT.500-14 for the assessment of TV pictures
\begin{itemize}
\item whose typical quality represents the extreme high levels of videotelephony and videoconferencing.
\end{itemize}
\item The specific comments of the DSIS scale, such as imperceptible or perceptible, are valuable when
\begin{itemize}
\item the viewer's detection of impairment is important.
\end{itemize}
\end{itemize}

\textbf{ACR} is easy and fast to implement, and the presentation of stimuli resembles ordinary system use. Therefore, ACR is well suited for qualification tests.

\textbf{ACR-HR} has all the advantages of ACR in terms of presentation and speed. Its main advantage over ACR is that the perceptual impact of the reference video can be removed from subjective scores. This reduces the influence of:

\begin{itemize}
\item \textbf{scene bias}, such as viewers liking or disliking a reference video;
\item \textbf{reference video quality}, such as small camera-quality differences;
\item \textbf{monitor} characteristics, such as professional versus consumer-grade displays.
\end{itemize}

ACR-HR is well suited to large experiments, provided that all reference videos are at least of "good" quality. However, ACR-HR may be insensitive to impairments that are easily detected by direct differential methods such as DSIS.

The main advantage of \textbf{PC} is its \textbf{high discriminatory power}, which is valuable when several test items are nearly equal in quality. If many items must be evaluated in one test, PC becomes lengthy. In that case, an ACR or DSIS test may first be carried out with a limited number of observers, followed by a PC test only on items that received similar ratings.

\subsection{Crowdsourcing}

Subjective quality assessment can be outsourced to workers from an internet crowd. Workers evaluate the quality of multimedia files in their own time, using their own hardware.
\begin{center}\includegraphics[width=0.600\linewidth]{Pasted_image_20260515113721.png}\end{center}
\textit{Figure 25: MTurk Marketplace scheme, where requesters have tasks that need to be completed and workers want to earn money and work on interesting tasks.}

Crowdsourcing challenges include:

\begin{itemize}
\item unknown displays;
\item unknown viewing distance;
\item different ambient lighting;
\item cheating participants;
\item network instability.
\end{itemize}

Collected crowdsourced scores require specific processing to remove unreliable voters.

\section{Objective Quality Evaluation}

Objective quality assessment algorithms predict visual quality by comparing a distorted signal with a reference, typically by modeling the human visual system.

Objective image quality assessment is based on mathematical models that predict perceived quality between a distorted image and a reference image. These methods consider characteristics of the human visual system to incorporate perceptual aspects in the quality measures. Modern methods increasingly rely on machine learning and large subjective datasets.

\subsection{Perceptual Effects of the Human Visual System}

Perceived brightness is not a simple function of light intensity.

In the \textbf{Mach band} pattern, the visual system tends to undershoot or overshoot around the boundaries of regions with different intensities.

\begin{center}\includegraphics[width=0.600\linewidth]{Pasted_image_20260515114134.png}\end{center}
\textit{Figure 27: Grayscale Mach band gradient comparing actual stepped illumination with perceived brightness, including overshoot and undershoot at transitions.}

In \textbf{simultaneous contrast}, a region's perceived brightness does not depend only on its intensity; it also depends on background intensity.

\begin{center}
\includegraphics[width=0.600\linewidth]{Pasted_image_20260515114152.png}\end{center}
\textit{Figure 28: Three identical gray squares on black, dark gray, and light gray backgrounds, where the same square appears lighter on a dark background and darker on a light background.}

Dynamic luminance effects show that the HVS can fill in non-existing information or wrongly perceive geometrical properties of objects.

\begin{center}
\includegraphics[width=0.600\linewidth]{Pasted_image_20260515114212.png}\end{center}
\textit{Figure 29: Rotating optical illusion made of black-and-white spiral sectors that appear to rotate.}

\begin{center}\includegraphics[width=0.600\linewidth]{Pasted_image_20260515114232.png}\end{center}
\textit{Figure 30: Hermann grid illusion, showing a black grid with bright corridors where dark illusory points appear at intersections.}

\begin{center}\includegraphics[width=0.600\linewidth]{Pasted_image_20260515114249.png}\end{center}
\textit{Figure 31: Anomalous motion illusion with a circular colored pattern whose blue, yellow, and green elements appear to rotate.}

Objective algorithms, or \textbf{metrics}, are desired when they estimate subjective MOS with high correlation.

Examples include:

\begin{itemize}
\item full-reference metrics: \textbf{PSNR}, \textbf{SSIM}, \textbf{VMAF};
\item no-reference metrics: \textbf{NIQE}, \textbf{BRISQUE}, \textbf{CLIP-IQA};
\item reduced-reference metrics: \textbf{RRED}.
\end{itemize}

\subsection{Full Reduced and No Reference Scenarios}

\begin{center}\includegraphics[width=0.600\linewidth]{Pasted_image_20260515114324.png}\end{center}
\textit{Figure 33: Full-reference block diagram where an input/reference signal passes through signal processing to produce an output/processed signal, and both the reference and output feed an FR metric.}

Objective metrics are categorized by reference availability:

\begin{itemize}
\item \textbf{Full-reference (FR)}: the original signal is available.
\item \textbf{Reduced-reference (RR)}: partial information about the original signal is available.
\item \textbf{No-reference (NR)}: the original signal is not available.
\end{itemize}

\subsection{Full-Reference Notation}

Full-reference methods depend on both the reference, or original, image and the test, or compressed, image.

Notation:

\begin{itemize}
\item $I(n,m,c,t)$ is the intensity of component $c$ at position $(n,m)$ and time $t$ in the reference
\begin{itemize}
\item signal.
\end{itemize}
\item $\hat{I}(n,m,c,t)$ is the corresponding value in the test signal.
\item The color space is typically YCbCr, although RGB can be used in some cases.
\end{itemize}

\subsection{MSE and PSNR}

\begin{calloutbox}{red!70!black}{Mean Square Error}
\textbf{Statement:} For a single image at time $t$ in a video sequence, the luminance-component \textbf{mean square error} is:

\[
\text{MSE}_Y(t) = \frac{1}{NM} \sum_{n,m}
\left(I(n,m,1,t) - \hat{I}(n,m,1,t)\right)^2
\]

\textbf{Interpretation:} MSE measures pixel-wise squared error, but it does not model perceptual relevance.
\end{calloutbox}

The logarithmic version of MSE is \textbf{PSNR}:

\begin{calloutbox}{red!70!black}{Peak Signal-to-Noise Ratio}
\textbf{Statement:} The luminance PSNR at time $t$ is:

\[
\text{PSNR}_Y(t) =
10\log_{10} \frac{255^2}{\text{MSE}_Y(t)}
\approx 48dB - 10\log_{10} \text{MSE}_Y(t)
\]

\textbf{Interpretation:} PSNR is easy to compute and widely used, but it measures signal fidelity more than perceived quality.
\end{calloutbox}

Typical PSNR interpretation:

\begin{itemize}
\item $\text{PSNR} > 45$ dB: excellent quality.
\item 40 to 45 dB: very good quality.
\item 30 to 40 dB: from good to very good quality.
\item below 30 dB: typically poor quality.
\end{itemize}

For chroma components, one can similarly define $\text{MSE}_{Cb}(t)$, $\text{MSE}_{Cr}(t)$, $\text{PSNR}_{Cb}(t)$, and $\text{PSNR}_{Cr}(t)$.

For example:

\[
\text{MSE}_{Cb}(t) =
\frac{1}{\frac{N}{2}\frac{M}{2}}
\sum_{n,m} \left(I(n,m,2,t) - \hat{I}(n,m,2,t)\right)^2
\]

A global quality measure for color images is:

\[
\text{PSNR}_{YCbCr}(t) =
\frac{3}{4}\text{PSNR}_Y(t)
+ \frac{1}{8}\text{PSNR}_{Cb}(t)
+ \frac{1}{8}\text{PSNR}_{Cr}(t)
\]

The smaller chroma weights reflect the lesser impact of color components on human perception.

For a video sequence, PSNR is averaged across the $T$ individual images:

\[
\text{PSNR}_Y = \frac{1}{T}\sum_{t=1}^{T} \text{PSNR}_Y(t)
\]

The same averaging applies to $\text{PSNR}_{Cb}$, $\text{PSNR}_{Cr}$, and $\text{PSNR}_{YCbCr}$.

\begin{center}\includegraphics[width=0.600\linewidth]{Pasted_image_20260515114654.png}\end{center}
\textit{Figure 34: Per-frame PSNR graph for the "foreman" sequence coded with H.264, distinguishing I, P, and B frames and reporting average $\text{PSNR}_Y \approx 39.0$ dB.}

\subsection{Bjontegaard Metrics}

Video codecs produce rate-distortion curves. A single-measure comparison can be obtained with \textbf{Bjontegaard metrics}.

\begin{enumerate}
\item \textbf{Bjontegaard Delta Rate (BD-Rate)} is the average rate difference at the same quality,
\begin{itemize}
\item expressed as a percentage. It is computed as the red area divided by the amplitude of the common
\item PSNR interval.
\end{itemize}
\item \textbf{Bjontegaard Delta PSNR (BD-PSNR)} is the average PSNR difference at the same rate, expressed
\begin{itemize}
\item in dB. It is computed as the blue area divided by the amplitude of the common rate interval.
\end{itemize}
\end{enumerate}

\begin{center}\includegraphics[width=0.600\linewidth]{Pasted_image_20260515114717.png}\end{center}
\begin{center}\includegraphics[width=0.600\linewidth]{Pasted_image_20260515114740.png}\end{center}

\textit{Figure 35: Rate-distortion curves for M-JP2K, MPEG-2, and H.264, with PSNR in dB plotted against rate in kbps and red and blue areas indicating BD-Rate and BD-PSNR computations.}

In practice, BD metrics are computed as follows:

\begin{enumerate}
\item For each curve, a small number of points, usually 4 or 5, is considered.
\item These points are fitted into a parametric representation of the curve, using a third-degree
\begin{itemize}
\item polynomial in the variable $t = \log R$.
\end{itemize}
\item The area is computed from the parametric representation of the two curves.
\end{enumerate}

\begin{center}\includegraphics[width=0.600\linewidth]{Pasted_image_20260515114955.png}\end{center}
\textit{Figure 36: Rate-distortion curve with the red area between H.264 and MPEG-2 highlighted for Bjontegaard computation.}

\subsection{Why PSNR Is Not Enough}

PSNR measures pixel fidelity, not perceptual quality.

Important limitations:

\begin{itemize}
\item all pixels are treated equally;
\item human visual perception is ignored;
\item different artefacts may produce similar PSNR;
\item temporal perception is poorly modeled.
\end{itemize}

\begin{calloutbox}{red!70!black}{PSNR Limitation}
\textbf{Statement:} Similar PSNR does not imply similar perceived quality.

\textbf{Interpretation:} Two videos or images may have nearly identical pixel errors but very different visual annoyance, such as blur versus low frame rate in the subjective test.
\end{calloutbox}

\subsection{SSIM Structural Similarity Index}

In 2004, the \textbf{Structural Similarity Index (SSIM)} was introduced as a new approach for image quality assessment.

\begin{itemize}
\item SSIM is a full-reference metric.
\item Its value lies in the interval $[0,1]$, from bad to good.
\item SSIM improves on traditional metrics such as PSNR and MSE, which are often inconsistent with
\begin{itemize}
\item human perception.
\end{itemize}
\item SSIM is based on the human visual system.
\end{itemize}
\begin{center}\includegraphics[width=0.600\linewidth]{Pasted_image_20260515115320.png}\end{center}
\textit{Figure 37: Reference "Barbara" image and three versions with additive Gaussian noise of increasing variance, 64, 200, and 600, with PSNR values 30.05 dB, 25.15 dB, and 20.42 dB, showing that PSNR decreases reliably when the same degradation increases.}

\begin{center}\includegraphics[width=0.600\linewidth]{Pasted_image_20260515115335.png}\end{center}
\textit{Figure 38: Reference "Barbara" image and three distorted versions with similar PSNR values, 20.42 dB, 20.42 dB, and 20.76 dB, corresponding to additive Gaussian noise, blur, and additive Gaussian noise in high frequencies, showing very different visual perception despite similar PSNR.}

Modern objective metrics attempt to model:

\begin{itemize}
\item human visual perception;
\item structural similarity;
\item temporal sensitivity;
\item perceptual importance.
\end{itemize}

The evolution is:

\[
\text{MSE} \rightarrow \text{PSNR} \rightarrow \text{SSIM} \rightarrow \text{VMAF}
\rightarrow \text{AI-based metrics}
\]

Modern metrics try to estimate \textbf{QoE}, not only signal fidelity.

The main idea of SSIM is that the human visual system is highly sensitive to:

\begin{itemize}
\item structures;
\item edges;
\item local contrast.
\end{itemize}

Instead of pure pixel error, SSIM compares:

\begin{itemize}
\item luminance;
\item contrast;
\item structural information.
\end{itemize}

\begin{center}
\begin{tabular}{|l|l|}
\hline
\textbf{SSIM} & \textbf{Quality} \\ \hline
1.0 & Identical \\ \hline
\textgreater{} 0.95 & Very high \\ \hline
0.8 to 0.9 & Noticeable degradation \\ \hline
\end{tabular}
\end{center}

\begin{calloutbox}{red!70!black}{Structural Similarity Index}
\textbf{Statement:} SSIM is defined as:

\[
SSIM(x,y) = [l(x,y)]^\alpha [c(x,y)]^\beta [s(x,y)]^\gamma
\]

\textbf{Interpretation:} SSIM combines luminance, contrast, and structure comparisons, with parameters $\alpha$, $\beta$, and $\gamma$ controlling their relative importance.
\end{calloutbox}

The components are:

\begin{itemize}
\item $l(x,y)$ accounts for luminance change and is qualitatively consistent with Weber's law, which
\begin{itemize}
\item models light adaptation in the HVS.
\end{itemize}
\item $c(x,y)$ accounts for contrast change and is consistent with contrast masking in the HVS.
\item $s(x,y)$ is the structure comparison function.
\end{itemize}

\begin{center}\includegraphics[width=0.600\linewidth]{Pasted_image_20260515115415.png}\end{center}
\textit{Figure 39: Image of colorful hats distorted by additive white Gaussian noise, with PSNR 28.8 and SSIM 0.48.}

\begin{center}\includegraphics[width=0.600\linewidth]{Pasted_image_20260515115437.png}\end{center}
\textit{Figure 40: Reference image of colorful hats.}

\begin{center}\includegraphics[width=0.600\linewidth]{Pasted_image_20260515115452.png}\end{center}
\textit{Figure 41: Image of colorful hats distorted by blur, with PSNR 33.78 and SSIM 0.91.}

\begin{center}\includegraphics[width=0.600\linewidth]{Pasted_image_20260515115505.png}\end{center}
\textit{Figure 42: Image of colorful hats distorted by JPEG compression, with PSNR 33.17 and SSIM 0.88.}

\begin{center}\includegraphics[width=0.600\linewidth]{Pasted_image_20260515115521.png}\end{center}
\textit{Figure 43: Image of colorful hats distorted by geometrical deformation, with PSNR 30.62 and SSIM 0.77.}

\subsection{VMAF}

\textbf{Video Multi-Method Assessment Fusion (VMAF)} is a full-reference video metric proposed by Netflix in June 2016.

VMAF combines multiple features:

\begin{itemize}
\item detail preservation;
\item structural similarity;
\item motion information;
\item perceptual learning.
\end{itemize}

VMAF values span from 0 to 100, where 0 is worst and 100 is best. A VMAF value of 93 is a good target, and content with an average over 90 usually looks good.

\begin{center}\includegraphics[width=0.600\linewidth]{Pasted_image_20260515115553.png}\end{center}
\textit{Figure 44: "How VMAF works" diagram, showing pixel neighborhoods feeding spatial feature extraction with VIF and DLM and temporal feature extraction with TI, then within-frame spatial pooling, SVM prediction or fusion, per-frame score, and temporal pooling. The model is trained using subjective data.}

\subsection{AI-Based Quality Metrics}

Modern quality metrics use:

\begin{itemize}
\item deep neural networks;
\item learned perceptual representations;
\item large subjective datasets.
\end{itemize}

Examples include:

\begin{itemize}
\item \textbf{LPIPS};
\item \textbf{DISTS};
\item learned no-reference metrics.
\end{itemize}

AI-based metrics can:

\begin{itemize}
\item model complex artefacts;
\item predict perceptual annoyance;
\item estimate quality without reference videos.
\end{itemize}

The bottom takeaway is that quality assessment is increasingly data-driven.

\subsection{Reduced-Reference Metrics}

\begin{center}\includegraphics[width=0.600\linewidth]{Pasted_image_20260515115621.png}\end{center}
\textit{Figure 45: Reduced-reference block diagram where the input/reference signal goes through signal processing to produce an output/processed signal, while features extracted from the reference and the processed output feed an RR metric.}

The key idea of \textbf{Reduced-Reference (RR)} metrics is that only partial information about the reference signal is available.

Typical transmitted features include:

\begin{itemize}
\item edge statistics;
\item frequency coefficients;
\item texture descriptors;
\item motion information.
\end{itemize}

Advantages:

\begin{itemize}
\item lower bandwidth requirements;
\item suitability for network monitoring;
\item practicality for large-scale systems.
\end{itemize}

Limitations:

\begin{itemize}
\item RR metrics are less accurate than full-reference metrics;
\item RR metrics require side-information transmission.
\end{itemize}

Typical applications include IPTV monitoring, adaptive streaming, and network QoE estimation.

\subsection{No-Reference Metrics}
\begin{center}\includegraphics[width=0.600\linewidth]{Pasted_image_20260515115703.png}\end{center}
\textit{Figure 46: No-reference block diagram where an input/reference signal passes through signal processing to produce an output/processed signal, and only the output/processed signal feeds an NR metric.}

The key idea of \textbf{No-Reference (NR)} metrics is that quality is estimated without access to the original signal.

NR metrics are important because, in many real systems:

\begin{itemize}
\item the reference video is unavailable;
\item only the received content can be analyzed.
\end{itemize}

Typical scenarios include:

\begin{itemize}
\item social media platforms;
\item smartphone cameras;
\item user-generated content;
\item video surveillance;
\item streaming analytics.
\end{itemize}

The main challenge is distinguishing:

\begin{itemize}
\item original scene characteristics;
\item distortion artefacts.
\end{itemize}

NR quality assessment is one of the most challenging problems in multimedia systems.

Examples of no-reference metrics:

\begin{center}
\begin{tabular}{|l|l|}
\hline
\textbf{Metric} & \textbf{Main idea} \\ \hline
BRISQUE & Natural scene statistics \\ \hline
NIQE & Deviation from natural image statistics \\ \hline
PIQE & Block-wise distortion estimation \\ \hline
Deep NR Metrics & Learned perceptual prediction \\ \hline
\end{tabular}
\end{center}

Modern AI-based no-reference quality assessment uses:

\begin{itemize}
\item convolutional neural networks;
\item transformers;
\item perceptual embeddings;
\item multimodal learning.
\end{itemize}

Goals:

\begin{itemize}
\item predict MOS directly;
\item estimate perceptual annoyance;
\item operate in real time.
\end{itemize}

Current trends:

\begin{itemize}
\item NR video quality for streaming;
\item mobile QoE prediction;
\item generative AI quality assessment;
\item perceptual metrics for AI-generated media.
\end{itemize}

AI is transforming quality assessment from signal analysis to perceptual prediction.

\section{Subjective Versus Objective Assessment}

Full-reference, reduced-reference, and no-reference methods can be compared as follows:

\begin{center}
\begin{tabular}{|l|l|l|l|l|}
\hline
\textbf{Type} & \textbf{Reference availability} & \textbf{Accuracy} & \textbf{Complexity} & \textbf{Typical use} \\ \hline
FR & Full reference & Very High & Medium & Codec evaluation \\ \hline
RR & Partial reference & Medium & Medium & Network monitoring \\ \hline
NR & No reference & Lower / evolving & High & Real-time QoE \\ \hline
\end{tabular}
\end{center}

Subjective and objective assessment have complementary properties:

\begin{center}
\begin{tabular}{|l|l|}
\hline
\textbf{Subjective} & \textbf{Objective} \\ \hline
\textbf{Accurate} & \textbf{Fast} \\ \hline
\textbf{Human perception} & \textbf{Automatic} \\ \hline
\textbf{Expensive} & \textbf{Scalable} \\ \hline
\textbf{Time-consuming} & \textbf{Real-time capable} \\ \hline
\textbf{Ground truth} & \textbf{Approximation} \\ \hline
\end{tabular}
\end{center}

\section{Final Takeaways}

The main content of the lecture is that multimedia quality is ultimately about human perception.

\begin{itemize}
\item \textbf{QoE} is fundamentally subjective.
\item Human observers do not always agree.
\item MOS requires statistical processing.
\item PSNR alone is insufficient.
\item Different artefacts affect perception differently.
\item Modern systems optimize perceptual quality.
\item AI-based metrics are becoming increasingly important.
\end{itemize}

\textgreater{} Multimedia quality is ultimately about human perception.

\section{Summary Table}

\begin{center}
\begin{tabular}{|l|l|l|l|}
\hline
\textbf{Concept} & \textbf{Definition / Formula} & \textbf{Compression role} & \textbf{Notes} \\ \hline
Quality assessment & Associates a quantitative grade with a multimedia stimulus & Provides the benefit term in cost-benefit analysis & Can be subjective or objective \\ \hline
Subjective assessment & Human observers directly rate quality & Gold standard for QoE & Accurate but expensive and time-consuming \\ \hline
Objective assessment & Algorithm predicts perceived quality & Enables automatic optimization and monitoring & Must correlate with subjective MOS \\ \hline
QoE & Quality as perceived by users & Final optimization target & Multidimensional and subjective \\ \hline
Coding rate & Number of bits used to represent/transmit content & Main cost in compression systems & Traded against quality \\ \hline
ACR & Single stimulus rated independently & Fast large-scale subjective testing & Does not directly test fidelity \\ \hline
ACR-HR & ACR with hidden reference & Removes some reference-content bias & Useful for large experiments \\ \hline
DSIS & Reference and impaired version are shown & Tests fidelity to the source & Accurate but slower than ACR \\ \hline
PC & Pairwise A/B preference comparison & Fine perceptual ranking & High discriminatory power but lengthy \\ \hline
MOS & $MOS = \frac{1}{N}\sum_{i=1}^{N}x_i$ & Aggregates subjective ratings & Needs variability and confidence analysis \\ \hline
Standard error & $SE = \frac{s}{\sqrt{N}}$ & Quantifies MOS uncertainty & Decreases with more observers \\ \hline
Confidence interval & $ci = [-1.96SE; 1.96SE]$ & Expresses statistical reliability & Based on repeated-experiment interpretation \\ \hline
SI & $SI = \max_{time}\{\sigma_{space}[Sobel(F_n)]\}$ & Characterizes spatial detail of video stimuli & High for textured scenes \\ \hline
TI & $TI = \max_{time}\{\sigma_{space}[M_n(i,j)]\}$ & Characterizes temporal activity & High for high-motion scenes \\ \hline
MSE & Average squared pixel error & Basis for PSNR & Easy but perceptually weak \\ \hline
PSNR & $10\log_{10}\frac{255^2}{\text{MSE}}$ & Common codec quality metric & Similar PSNR can mean different perceived quality \\ \hline
YCbCr-PSNR & Weighted average of Y, Cb, and Cr PSNR & Color-image quality score & Luma receives larger weight \\ \hline
BD-Rate & Average rate difference at the same quality & Compares rate-distortion curves & Expressed as percentage \\ \hline
BD-PSNR & Average PSNR difference at the same rate & Compares rate-distortion curves & Expressed in dB \\ \hline
SSIM & $[l(x,y)]^\alpha[c(x,y)]^\beta[s(x,y)]^\gamma$ & Perceptual full-reference metric & Compares luminance, contrast, and structure \\ \hline
VMAF & Learned fusion of perceptual video features & Practical video quality prediction & Combines spatial, temporal, and learned components \\ \hline
FR metric & Uses full reference signal & Codec evaluation & Highest reference availability \\ \hline
RR metric & Uses partial reference features & Network monitoring & Requires side information \\ \hline
NR metric & Uses no reference signal & Real-time QoE estimation & Hardest objective assessment setting \\ \hline
AI-based metrics & Learned perceptual prediction from data & Model complex artefacts & Increasingly important for modern media \\ \hline
\end{tabular}
\end{center}

